\documentclass[8pt,a4paper]{extarticle}

\usepackage[margin=0.42in,top=0.50in,bottom=0.46in]{geometry}
\usepackage{amsmath,amssymb}
\usepackage{array,booktabs}
\usepackage{enumitem}
\usepackage{multicol}
\usepackage{xcolor}
\usepackage{fancyhdr}
\usepackage{hyperref}
\hypersetup{colorlinks=true,urlcolor=DeepBlue}

\definecolor{DeepBlue}{HTML}{173B57}
\definecolor{MidBlue}{HTML}{2D6A8A}
\definecolor{LightBlue}{HTML}{EAF3F8}
\definecolor{Warm}{HTML}{FFF4DD}
\definecolor{Ink}{HTML}{17212B}
\definecolor{Muted}{HTML}{596773}

\pagestyle{fancy}
\fancyhf{}
\lhead{\scriptsize\color{Muted}Scaling Book - Chapter 7}
\rhead{\scriptsize\color{Muted}Transformer Inference}
\cfoot{\scriptsize\color{Muted}\thepage}
\renewcommand{\headrulewidth}{0.3pt}
\setlength{\headheight}{12pt}

\setlength{\parindent}{0pt}
\setlength{\parskip}{1.1pt}
\setlength{\columnsep}{12pt}
\setlength{\columnseprule}{0.25pt}
\setlength{\multicolsep}{2pt}
\setlist[itemize]{leftmargin=1.2em,itemsep=0.2pt,topsep=0.8pt}
\setlist[enumerate]{leftmargin=1.4em,itemsep=0.2pt,topsep=0.8pt}
\setlength{\abovedisplayskip}{1.5pt}
\setlength{\belowdisplayskip}{1.5pt}
\setlength{\abovedisplayshortskip}{1pt}
\setlength{\belowdisplayshortskip}{1pt}

\newcommand{\sect}[1]{%
  \vspace{2pt}%
  {\color{DeepBlue}\normalsize\bfseries #1}\par
  \vspace{0.5pt}\hrule\vspace{1.5pt}
}
\newcommand{\subsect}[1]{%
  \vspace{1pt}{\color{MidBlue}\bfseries #1}\par
}
\newcommand{\formula}[1]{%
  \[
    \boxed{\displaystyle #1}
  \]
}
\newcommand{\callout}[2][LightBlue]{%
  \vspace{1pt}\noindent
  \colorbox{#1}{%
    \parbox{\dimexpr\linewidth-2\fboxsep\relax}{\footnotesize #2}%
  }\vspace{1pt}
}

\begin{document}

\begin{center}
{\color{DeepBlue}\Large\bfseries Transformer Inference}\\[-1pt]
{\normalsize Worked Problems - Condensed Review}\\[1pt]
{\scriptsize Lucas Yan \quad$\cdot$\quad
\href{https://jax-ml.github.io/scaling-book/inference/}{JAX Scaling Book, Chapter 7}}
\end{center}

\callout{%
\textbf{Central idea.} Decode is usually a \textbf{data-movement problem}, not
a FLOPs problem. Start from tensor shapes, convert to bytes, convert bytes to
time, then identify the largest bottleneck.
}

\begin{multicols}{2}

\sect{1. Symbols that must stay separate}

\begin{tabular}{@{}>{$}l<{$} p{0.72\linewidth}@{}}
\toprule
B & concurrent decode requests / tokens\\
T & tokens processed now; decode usually $T=1$\\
S & cached context length\\
L,D,F & layers, model width, MLP width\\
N,K,H & query heads, KV heads, head dimension\\
E,k & total experts, activated experts/token\\
b_w,b_{\mathrm{KV}} & bytes per weight / KV element\\
\bottomrule
\end{tabular}

\callout[Warm]{%
\textbf{Do not mix these:}
$N$ makes queries; $K$ determines cached heads.
$E$ experts are stored; only $k$ are used by one token.
The $2$ in KV memory means \emph{key and value}, not $k=2$ experts.
}

\sect{2. Dense-model parameter count}

For a gated MLP and shared input/output embedding:
\formula{P=3LDF+2LDH(N+K)+DV}

\[
P_{\mathrm{MLP}}=3LDF,\qquad
P_{\mathrm{attn}}=2LDH(N+K).
\]

Attention decomposition:
\[
\underbrace{W_Q+W_O}_{2DNH}
\;+\;
\underbrace{W_K+W_V}_{2DKH}.
\]

\textbf{Worked model:}
\[
P=12.885\mathrm{B}+5.369\mathrm{B}+0.132\mathrm{B}
\approx\boxed{18.4\mathrm{B}}.
\]

Parameters are counts, not bytes:
\formula{M_{\mathrm{weights}}=b_wP}
Int8 uses 1 byte/parameter; bf16 uses 2.

\sect{3. KV cache}

Only keys and values survive for future tokens:
\formula{m_{\mathrm{KV/token}}=2LKH\,b_{\mathrm{KV}}}

For one sequence:
\formula{M_{\mathrm{KV,seq}}=S(2LKH)b_{\mathrm{KV}}}

Conceptual shape:
\[
\mathrm{KV}[L,2,B,S,K,H].
\]

\textbf{Worked model, int8:}
\[
2(64)(8)(256)=262{,}144\text{ bytes/token}
\approx\boxed{262\text{ kB/token}}.
\]

At $S=128{,}000$:
\[
M_{\mathrm{KV,seq}}\approx\boxed{33.55\text{ GB}}.
\]

\sect{4. Memory-capacity questions}

Weights are a fixed cost; every request adds another KV cache:
\formula{
B_{\max}=
\left\lfloor
\frac{N_cM_{\mathrm{HBM}}-M_{\mathrm{weights}}}
{M_{\mathrm{KV,seq}}}
\right\rfloor
}

For 16 TPU v5e chips:
\[
B_{\max}=
\left\lfloor\frac{256-18.4}{33.55}\right\rfloor
=\boxed{7}.
\]

Reducing $K:8\to1$ makes KV memory $8\times$ smaller:
\[
\boxed{B_{\max}\approx56}.
\]

\sect{5. Bytes become latency}

Universal conversion:
\formula{\text{time}=\frac{\text{bytes moved}}{\text{bandwidth}}}

Fully sharded weight loading:
\formula{
T_{\mathrm{weights}}=
\frac{b_wP}{N_cW_{\mathrm{HBM}}}
}

Worked result:
\[
\frac{18.4\times10^9}{16(8.2\times10^{11})}
\approx\boxed{1.4\text{ ms}}.
\]

All chips load their shards simultaneously; use aggregate bandwidth
$N_cW_{\mathrm{HBM}}$.

\sect{6. Prefill versus decode}

\begin{tabular}{@{}p{0.25\linewidth}p{0.31\linewidth}p{0.31\linewidth}@{}}
\toprule
 & \textbf{Prefill} & \textbf{Decode}\\
\midrule
Tokens now & many & one/request\\
Weight reuse & high & low\\
KV action & create & read history\\
Typical limit & compute & HBM bandwidth\\
Extra scaling & sequence & model/KV\\
\bottomrule
\end{tabular}

\callout{%
Each decode step produces one new token but reads $S$ previous KV positions.
Long context therefore makes attention bandwidth-heavy.
}

\sect{7. Decode-step latency}

\[
T_{\mathrm{KV}}=
\frac{BM_{\mathrm{KV,seq}}}{N_cW_{\mathrm{HBM}}},
\qquad
T_{\mathrm{math}}=
\frac{2BP_{\mathrm{act}}}{N_cC}.
\]

Core lower bound:
\formula{
T_{\mathrm{step}}\approx
T_{\mathrm{KV}}+
\max(T_{\mathrm{weights}},T_{\mathrm{math}})
}

Why add then max?
\begin{itemize}
  \item Attention and linear blocks are sequential, so their times add.
  \item Inside a linear block, compute and weight loading compete; slower wins.
\end{itemize}

At $B=7,S=128{,}000$:
\[
T_{\mathrm{KV}}\approx17.9\text{ ms},\quad
T_{\mathrm{weights}}\approx1.4\text{ ms},\quad
T_{\mathrm{math}}\approx0.041\text{ ms}.
\]
\[
\boxed{T_{\mathrm{step}}\approx19.3\text{ ms}}
\]
before fixed collective/runtime overheads.

\sect{8. Rooflines}

A roofline is a \textbf{crossing of two candidate bottlenecks}.

\begin{itemize}
  \item ``When FLOPs-bound?'' Set $T_{\mathrm{math}}=T_{\mathrm{HBM}}$.
  \item ``How far can tensor parallelism scale?'' Set
  $T_{\mathrm{ICI}}=T_{\mathrm{HBM}}$.
\end{itemize}

Decode tensor-parallel limit:
\formula{
p_{\max}\approx\frac{F}{B\beta},
\qquad
\beta=\frac{W_{\mathrm{HBM}}}{W_{\mathrm{ICI}}}
}

Prefill approximation:
\formula{p_{\mathrm{prefill}}\approx F/2200}

For $F=16{,}384$, $B=7$, $\beta\approx8$:
\[
p_{\max}^{\mathrm{decode}}\approx293\gg16,
\qquad
p_{\mathrm{prefill}}\approx7.4.
\]

Interpretation: 16-way decode remains HBM-limited in the simplified model;
prefill should use roughly 4--8-way tensor parallelism, then sequence
parallelism.

\sect{9. KV-cache sharding}

Given $\mathrm{KV}[2,B,S,K,H]$, consider:
\begin{enumerate}
  \item \textbf{$K$ first}: different devices own KV heads.
  \item \textbf{$B$ second}: devices own different requests.
  \item \textbf{$S$ third}: devices own context segments.
  \item Avoid replication when the cache is large.
\end{enumerate}

Head sharding is limited to $K$ ways. Beyond that, batch/sequence sharding
adds collectives, commonly AllToAlls for layout changes.

\callout[Warm]{%
\textbf{Do not invent communication bytes} by adding $BD+DF+BF$.
Count only tensors actually transmitted by each collective. Stationary
weights remain in HBM; activations usually move over ICI.
}

\end{multicols}

\newpage

\begin{center}
{\color{DeepBlue}\large\bfseries MoE and Large-Scale Sharding}
\end{center}

\begin{multicols}{2}

\sect{10. MoE: stored versus activated}

Total stored parameters:
\formula{
P_{\mathrm{total}}=
3ELDF+2LDH(N+K)+DV
}

Used by one token:
\formula{
P_{\mathrm{activated}}=
3kLDF+2LDH(N+K)+DV
}

Forward FLOPs for $T$ tokens:
\formula{\mathrm{FLOPs}\approx2P_{\mathrm{activated}}T}

Worked values for $E=16,k=2$:
\[
P_{\mathrm{total}}\approx\boxed{212\mathrm{B}},
\qquad
P_{\mathrm{activated}}\approx\boxed{31.2\mathrm{B}},
\]
\[
\mathrm{FLOPs}\approx\boxed{62.4\times10^9T}.
\]

MoE adds $E$ times the expert storage but only $k$ times expert compute.
Therefore:
\formula{
B_{\mathrm{crit,MoE}}
\approx B_{\mathrm{crit,dense}}\frac{E}{k}
}

\[
240\frac{16}{2}=\boxed{1920\text{ tokens}}.
\]

\callout{%
\textbf{MoE changes the MLP, not attention.}
The KV-cache formula $2LKH$ is unchanged.
}

\sect{11. Expert sharding}

Expert weight tensor:
\[
W[E,D,F]\longrightarrow W[E_Z,D_X,F_Y].
\]

\begin{itemize}
  \item $Z$: which experts a chip owns.
  \item $Y$: which $F$ columns of that expert.
  \item $X$: which $D$ rows; training/FSDP dimension here.
\end{itemize}

At inference, $X=1$. Idealized weight balance:
\formula{
M_{\mathrm{weights/chip}}\approx\frac{b_wP}{YZ}
}
\[
T_{\mathrm{HBM}}\approx
\frac{b_wP}{YZ\,W_{\mathrm{HBM}}},
\qquad
M_{\mathrm{free}}\approx
M_{\mathrm{chip}}-\frac{b_wP}{YZ}.
\]

Minimum capacity:
\formula{
N_{\min}=
\left\lceil\frac{b_wP}{M_{\mathrm{chip}}}\right\rceil
}
then round up to a supported, divisible topology.

For the 212B MoE in bf16:
\begin{itemize}
  \item $Y=8,Z=16$ ($128$ chips): $\approx4.0$ ms weight load,
  $\approx12.7$ GB free/chip.
  \item Smallest practical slice: $\boxed{32\text{ chips}}$.
\end{itemize}

\callout[Warm]{%
\textbf{Memory rule:} experts stay; tokens travel.
Expert parallelism avoids moving huge weights by routing much smaller
activations to expert-owning chips.
}

\columnbreak

\sect{12. 2D weight-stationary sharding}

Split both $D$ and $F$ so local weight tiles become more balanced.
Let total chips be $N=XYZ$.

Two local matmuls:
\formula{
T_{\mathrm{math}}=\frac{4BDF}{NC}
}

Collective time:
\formula{
T_{\mathrm{comms}}=
\frac{2BD}{XW_{\mathrm{ICI}}}
+
\frac{4BF}{YZ\,W_{\mathrm{ICI}}}
}

For $F=4D$:
\[
T_{\mathrm{comms}}
=
\frac{2BD}{W_{\mathrm{ICI}}}
\left(\frac1X+\frac8{YZ}\right).
\]

At the optimum, balance the communication terms:
\[
\frac1X=\frac8{YZ}
\quad\Longrightarrow\quad
\boxed{YZ=8X}.
\]

Using $X(YZ)=N$:
\[
\boxed{
X=\sqrt{\frac N8},
\qquad
YZ=\sqrt{8N}
}.
\]

Optimal communication:
\[
\boxed{
T_{\mathrm{comms}}^{2D}
\approx
\frac{\sqrt{128}\,BD}
{\sqrt N\,W_{\mathrm{ICI}}}
}.
\]

It beats the chapter's traditional baseline when:
\[
\boxed{N>72}
\qquad
\text{or generally}\qquad
\boxed{N>18F/D}.
\]

\textbf{Meaning:} 2D sharding adds collectives, but makes the communicated
buffers shrink with scale. It pays off only after enough chips.

\sect{13. Recurring traps}

\begin{enumerate}
  \item \textbf{Parameters vs bytes:} multiply by bytes/parameter.
  \item \textbf{Total vs activated:} storage uses $E$; token FLOPs use $k$.
  \item \textbf{Query vs KV heads:} $N$ affects query projection; $K$ affects
  cache size.
  \item \textbf{Wrong roofline pair:} choose math/HBM or ICI/HBM based on the
  question.
  \item \textbf{Wrong aggregation:} fully sharded chips load simultaneously;
  both bandwidth and compute scale with chip count.
  \item \textbf{Wrong max/sum:} sequential stages add; competing bottlenecks
  use the maximum.
\end{enumerate}

\sect{14. Seven-step solving checklist}

\begin{enumerate}
  \item Write the tensor shape.
  \item Count elements.
  \item Multiply by bytes/element.
  \item Mark cost as fixed, per token, or per request.
  \item Apply sharding dimensions.
  \item Convert work/bytes into candidate times.
  \item Name the largest bottleneck and sanity-check units.
\end{enumerate}

\callout{%
\textbf{Mastery test.}
If $K$ halves, $E$ doubles, and $k$ stays fixed:
KV memory halves; total expert parameters roughly double; activated expert
parameters stay fixed; the MoE critical batch roughly doubles.
}

\vfill
{\small\color{Muted}
Reference: \url{https://jax-ml.github.io/scaling-book/inference/}\\
Summary derived from Worked Problems 1--7 and the accompanying tutoring
discussion.}

\end{multicols}

\end{document}
